\documentclass[11pt]{article}
\usepackage[a4paper,margin=25mm]{geometry}
\usepackage{amsmath,amssymb,amsthm,booktabs}
\newtheorem{theorem}{Theorem}\newtheorem{lemma}{Lemma}
\newcommand{\Aut}{\operatorname{Aut}}\newcommand{\HS}{\mathrm{HS}}
\title{Local data of the Higman--Sims coset geometries}
\author{}\date{}
\begin{document}\maketitle
We acquire the reviewed classification of 86 coset geometries for the
literal faithful group $G\cong\HS$, of order $44\,352\,000$.
The numbers in ranks two through six are $6,30,41,8,1$.
The acquired action is flag-transitive. Its parabolics are denoted by
$G_i$, $i\in I$, and $G_J=\bigcap_{j\in J}G_j$, with $G_\varnothing=G$.
The Borel subgroup is $B=G_I$. All identifiers below refer to the fixed
representatives in \texttt{full\_fusion03/result.g}.

\begin{theorem}
The retained finite calculations determine, for each of these 86
representatives, its type sizes, every flag stabilizer order, its Borel
and panel orders, its node orders, and the orders and literal generators
of its local action kernels. They also determine the two ordered
diameters and girth of each of its 437 rank-two residues at a base flag.
Finally, they determine the normalizer of the specified $G$ action in
the type-preserving automorphism group. For representatives 14 and 15
this is the full automorphism group, of respective orders
$44\,352\,000$ and $88\,704\,000$.
\end{theorem}

The assertion is a computation of the actual models. It does not assume
agreement with any printed group name or diagram label. The historical
supplement of 6 October 2003, attributed to Dimitri Leemans, is compared
with these values in a separate source audit. No new classification or
open-conjecture resolution is asserted here.

\section*{Stabilizers and local kernels}
The objects of type $i$ are the right cosets of $G_i$, so their number
is $|G:G_i|$. The stabilizer of the standard flag of type $J$ is exactly
$G_J$. Flag transitivity gives $|G:G_J|$ flags of that type. Fixing the
standard flag of cotype $\{i\}$ leaves $|G_{I\setminus\{i\}}:B|$
objects; consequently the node order is this index minus one.

Fix an object of type $i$. For $j\ne i$, its incident objects of type
$j$ form the transitive $G_i$-set with stabilizer $G_i\cap G_j$.
The kernel on that set is the core of this intersection in $G_i$.
An element acts trivially on the whole residue precisely when it acts
trivially on every type. Thus the local kernel is
\[
 K_i=\bigcap_{j\ne i}\operatorname{Core}_{G_i}(G_i\cap G_j).
\]
The calculation uses these actual intersections and cores. It saves
generators of $B$, the panels, each $K_i$, and the original parabolics,
as well as their orders. An order by itself is not a structural group
identification.

\section*{An indexed construction of the incidence graphs}
For a residue of cotype $\{i,j\}$ put
$H=G_{I\setminus\{i,j\}}$, $P=H\cap G_i$, $Q=H\cap G_j$,
and $D=P\cap Q=B$. Its two types are the cosets of $P$ and $Q$ in $H$.
Take right transversals $T_P,T_Q$ in $H$ and a right transversal $T_D$
of $D$ in $P$. For $x\in T_P$, its full neighbour set is
\begin{equation}\label{neighbours}
                     \{Qux:u\in T_D\}.
\end{equation}
Indeed a coset meeting $Px$ contains some $px$, $p\in P$, and is
therefore $Qpx$. Two such cosets agree precisely when
$p_1p_2^{-1}\in P\cap Q=D$. Thus (\ref{neighbours}) contains each
neighbour exactly once. A position lookup in $T_Q$ gives the integer
vertex label without constructing a large permutation action.

For the six rank-two geometries this construction produces 158,447
vertices and 537,835 edges in total. The completed invocation took
34.36 seconds. An earlier generic action construction reached its
120-second bound without returning a graph; that record is preserved.
The new construction changes the indexing method, not the incidence
relation. Independent small controls compare every edge with literal
coset intersection.

For higher ranks, 425 graphs were obtained by transporting the base
$P$-orbit in the action of $H$ on $Q$-cosets. This is the same
neighbour formula. Six graphs for representatives 14 and 15 are
acquired from their earlier reviewed comparison. Together the 437
graphs have 187,005 vertices and 601,049 edges. The exact roster
contains one graph for every unordered type pair of every representative.
The reader checks every adjacency, both valencies, the two flag indices,
and $|H:B|$ edges. These checks supplement the exact coset construction;
they do not substitute matching counts for incidence.

\section*{Distances and the diagram parameters}
\begin{lemma}
In a finite connected bipartite graph containing a cycle, whose automorphism group is
transitive on each part, breadth-first search at one vertex in each
part determines both type diameters and the girth.
\end{lemma}
\begin{proof}
The eccentricity is constant on each part by transitivity. Hence the
largest distance found from its chosen vertex is that type's diameter.
For a breadth-first tree rooted at $v$, a non-tree edge $xy$ gives the
number $d(v,x)+d(v,y)+1$. Its tree paths and the edge contain a cycle
of length at most this number, so every such number is at least the
girth. Choose a shortest cycle and move a vertex of it to $v$ using
transitivity on the appropriate part. Its length is even, say $2m$.
There is no shorter path between vertices of this cycle than the
shorter arc, since such a path would give a shorter cycle. The vertex
opposite $v$ therefore has distance $m$, and its two cycle neighbours
have distance $m-1$. At least one of these two edges is not a tree
edge and gives $2m$. The minimum recorded number is consequently the
girth. This argument applies to either part.
\end{proof}
The acquired coset action is transitive on each part. All valencies
are at least two, so these finite graphs contain cycles and the lemma
applies. The two searches must agree on the girth.
The diagram triple is the left type diameter, half the girth, and the
right type diameter. A digon is checked as a complete bipartite graph.
A complete-graph residue is recognized only when one type consists
of all distinct two-subsets of the other, with exactly that incidence.

\section*{Normalizers and the two remaining ambiguous profiles}
Let $N\cong\Aut(G)$ be the acquired faithful automorphism action on $G$,
and put $S=\bigcap_i N_N(G_i)$. Let $M$ be the normalizer of the
specified $G$ action in the type-preserving geometry automorphism group.
The chamber stabilizer $M_\Phi$ maps by conjugation into $S$.
This map is injective: an element fixing $\Phi$ and commuting with $G$
fixes every chamber, because $G$ is chamber-transitive, and hence fixes
every object. Conversely each automorphism in $S$ induces the
chamber-fixing coset map $G_i x\mapsto G_i\alpha(x)$. Thus
$M_\Phi\cong S$. Since $G\le M$ is chamber-transitive,
\[
                         |M|=|G:B|\,|S|.
\]
This formula computes a normalizer; it does not in general prove that
all geometry automorphisms normalize the given copy of $G$.

For representatives 14 and 15, that last assertion is already paid by
the reviewed shadow argument. Each defines a graph on its 1100-point
type; restriction to that type is faithful, and the derived subgroup
of the full shadow-graph group is precisely the specified $G$.
Every geometry automorphism therefore normalizes $G$.
Their Borel orders are both 24, while the calculated orders of $S$
are respectively 24 and 48. The displayed formula gives the two full
orders stated above. All three type sizes are different, so their
correlation groups preserve types as well. For representative 15 the
resulting group is the full shadow-graph group, which is isomorphic
to $\Aut(G)$ by its faithful conjugation action. Representative 14
has full group equal to the specified $G$ itself. To justify the faithful
conjugation assertion, the calculation also checks $S\cap G=B$ in both
cases. An automorphism commuting with $G$ takes the base chamber to
$\Phi^g$ for some $g\in G$; equality of the object stabilizers forces
$g\in S\cap G=B$. It therefore fixes the base chamber and, by
transitivity, the entire geometry. The conjugation kernel is trivial.

\paragraph{Attribution and computational boundary.}
The group model, complete family classification, full automorphism
action and faithful shadow argument are acquired reviewed results.
The general arguments here are elementary coset and group-action
facts; their explicit derivations, numerical applications and validation
programs are supplied here. Exact GAP permutation-group algorithms
remain computational premises. No historical source label, full
automorphism claim for the other 84 models, abstract identification of
local kernels, Lean verification, measurement or whole-manuscript
completion follows from this component alone.
\end{document}
